We identify and solve a fundamental limitation of phase-field Allen-Cahn lattice Boltzmann method (LBM) for simulating droplet dynamics on curved surfaces: the Allen-Cahn sharpening term resists geometric wetting boundary condition corrections, suppressing contact angle accuracy on curved hydrophobic surfaces. We propose a geometric amplification factor $\alpha_{geo}$ that overcomes this resistance by over-correcting the interface gradient near the wall. Through systematic calibration across three parameter dimensions---contact angle ($\theta_{eq} = 90°$--$162°$), curvature ratio ($R^* = 0.5$--$3.0$), and grid resolution ($N = 80$--$180$)---we establish three key thresholds: (1) $\alpha_{geo}$ is effective only for $\theta_{eq} > 120°$ (strong hydrophobicity), (2) $\alpha_{geo}$ is most effective for small $R^*$ (high curvature), and (3) $\alpha_{geo}$ requires $D_0/\xi > 3.0$ (sufficient interface resolution). The optimal $\alpha_{geo} = 1.5$ is recommended as a universal default for superhydrophobic curved surfaces. Using this correction, we construct the We $\times$ $R^*$ phase diagram of asymmetric droplet spreading via Allen-Cahn LBM at a density ratio of 828:1, revealing three regimes: symmetric ($k < 1.2$), moderately asymmetric ($1.2 \leq k < 2.0$), and highly asymmetric ($k \geq 2.0$). The asymmetry ratio $k$ peaks at $\mathrm{We} \approx 15$ ($k = 3.30$ at $R^*=1.0$) and decreases at higher Weber numbers, consistent with experimental observations of droplet breakup. Validation across six curvature ratios and four contact angles shows quantitative agreement within $2.1\%$ of the experimental benchmark, while concave geometries correctly yield symmetric spreading ($k=1.0$). These results demonstrate that geometric amplification is essential for Allen-Cahn LBM to accurately capture wetting physics on curved surfaces.
